\section{Each Element in Its Own Setting}
\label{sec:each-element}

Each part of the formulary belongs first where it comes from: a psalm to its
psalm, a lesson to its letter, an oration to the books that first carried it.
That setting follows for all ten elements in the order of the Mass, with the
textual facts and the principal witnesses on which the three readings draw.
Each witness's argument is developed most fully in the reading that relies on
it.

\subsection{Introit: the just God and the servant's plea}

The Introit is taken from Ps~118, the long alphabetical psalm in praise of
God's law. Its antiphon joins two verses from different strophes. Verse 137
opens the strophe \latin{Sade} (vv.~137--144), whose theme is the justice of
God and of his testimonies: ``Thou art just, O Lord: and thy judgment is
right.'' Verse 124 belongs to the strophe \latin{Ain} (vv.~121--128), in which
the servant asks for help and instruction: ``Deal with thy servant according to
thy mercy.'' The chant sings only this first half and leaves aside ``and teach
me thy justifications''. The confession of justice and the plea for mercy thus
stand side by side with nothing between them. The psalm verse is the psalm's
opening beatitude, ``Blessed are the undefiled in the way''. Two verses before
the antiphon's opening verse~137, in the strophe the antiphon passes over
between verses 124 and 137, the psalm prays
\latin{Faciem tuam illumina super servum tuum} (v.~135), ``Make thy face to
shine upon thy servant'', with the verb and construction the Offertory will use
of the sanctuary.

Ildefonso Schuster notes that Ps~118 was used in Rome as a processional
psalm, sung on the afternoon of Good Friday as the procession passed from the
Lateran to the Sessorian basilica. Augustine, Hilary of Poitiers and
Cassiodorus hear the two verses as the words of a penitent who asks mercy and
not reward, and St Robert Bellarmine reads the \latin{iudicium} of verse 137 as
God's law; their reasoning carries the third reading.

\subsection{Collect: to shun the devil's contagion and follow God alone}

The Collect asks one gift under two infinitives: that God's people may shun
\latin{diabolica contagia}, the defilements that come from contact with the
devil, and follow God alone, \latin{te solum Deum}, \latin{pura mente}. The verb
\latin{sectari}, a frequentative of \latin{sequi}, means to follow after
closely, to pursue. The first petition turns the heart from what corrupts it,
and the second turns it to the one it is to seek. The 1861 English renders the
second ``with a clean heart follow thee, the only true God''.

The prayer's own history bears on its wording. Its earliest witness, the Old
Gelasian Sacramentary (Vat.\ Reg.\ lat.~316, in H.~A. Wilson's edition, Liber
III, sect.~XIII), places it in a Sunday Mass without number and reads \latin{te
solum Dominum puro corde sectari}. The eighth-century Frankish Gelasians of
Rheinau and St Gall set it, with this Sunday's Secret and Postcommunion, at the
twenty-first week after Pentecost. The Supplement to the Gregorian
Sacramentary sets it at the eighteenth Sunday after Pentecost, where one
manuscript reads \latin{dominum} and another \latin{deum}, both with \latin{pura
mente}. The 1962 wording \latin{te solum Deum pura mente} is thus the later
form of the prayer. Its \latin{mente} meets the Gospel's \latin{in tota mente
tua}, as the older \latin{corde} would have met \latin{ex toto corde tuo}. The
1962 Missal borrows the second clause once more, in a postcommunion that names
St Jane Frances: \latin{fácias terréna despícere, et te solum Deum pura mente
sectári}. There it asks for a heart that despises earthly things and follows
God alone.

\subsection{Epistle: a life worthy of one calling}

The lesson opens the moral part of the Letter to the Ephesians (chapters
4--6). In the Vulgate it begins \latin{Obsecro itaque vos}, ``I therefore
beseech you'', drawing the consequence of the first three chapters. Those
chapters set out the mystery of the Church as Christ's body, in which Jews and
gentiles are made one: ``he is our peace, who hath made both one \dots\ by him
we have access both in one Spirit to the Father'' (Eph 2:14, 18). They end
with Paul's prayer and doxology (3:14--21). The Missal's \latin{Fratres}
replaces the ``therefore'', and its closing \latin{Qui est benedictus in
saecula saeculorum. Amen} gives the lesson a doxology of its own. Verses 1--3
describe the manner of life the calling asks: humility, mildness and
patience, mutual support in charity, and care for ``the unity of the Spirit in
the bond of peace''. Verses 4--6 give its ground in a sevenfold ``one'': one
body, one Spirit, one hope, one Lord, one faith, one baptism, one God and
Father of all. Verse 7, which the Mass does not read, turns to the diverse
gifts given to each.

The lesson stands in the Missal's course of Pauline epistles for the Sundays
after Pentecost. The Sixteenth Sunday reads Eph 3:13--21. After 1~Cor 1:4--8
on the Eighteenth, Ephesians returns on the Nineteenth, Twentieth and
Twenty-first Sundays (4:23--28; 5:15--21; 6:10--17). Its principal
interpreters, St John Chrysostom, St Jerome, Theodoret of Cyrus and St Thomas
Aquinas, read verses 1--3 as a school of humility and forbearance; at verse 6
they follow two lines, set out in the second reading.

\subsection{Gradual: the blessed nation and the creating word}

The Gradual takes two verses of Ps~32, a hymn whose Vulgate title is
\latin{Psalmus David}. Its first part praises the word and counsel of God that
created all things and stand for ever (vv.~4--11). Its centre blesses the
nation he chose (v.~12). Its close watches over those who fear him
(vv.~13--22). The Gradual reverses the psalm's order: the respond is the
beatitude of verse 12 and the versicle the creation verse 6. Its wording,
\latin{Dominus Deus eorum \dots\ quem elegit Dominus}, is not the Clementine
Vulgate's, but it is the reading of the psalter Augustine expounded
(\work{Enarrationes in Psalmos} 32, I, 12).

The same Gradual serves three other Masses of the 1962 Missal: the Wednesday
of the fourth week of Lent, where it follows the lesson ``Wash yourselves, be
clean'' (Isa 1:16--19); the Wednesday of the September Ember Days, three days
after this Sunday; and the votive Mass of the Holy Spirit. In the votive Mass
the versicle is printed with a capital, \latin{Spíritu oris eius}; at this
Sunday it is printed \latin{spíritu}. Schuster observes of this Sunday that, ``whether by coincidence or
design, it carries on the idea which the Epistle began to expound''. Its verse 12, the people God chose for his inheritance, belongs chiefly to
the first reading; its verse 6, the heavens made firm by the word and the
spirit of the Lord's mouth, to the second.

\subsection{Alleluia: the cry of the poor man}

The Alleluia verse is the first petition of Ps~101, whose title the Douay
renders ``The prayer of the poor man, when he was anxious, and poured out his
supplication before the Lord''. The psalm moves from the afflicted man's
lament (vv.~2--12) to God's enduring reign and the rebuilding of Sion (``Thou
shalt arise and have mercy on Sion'', v.~14), and to God's permanence against
the heavens that ``shall perish but thou remainest'' (vv.~26--28). It is the fifth of the seven penitential
psalms. Bellarmine notes that ``This verse is used daily by the Church as a
preparation to any other petitions''; Schuster, that it precedes almost all the
solemn prayers of the Church. The Missal's \latin{perveniat} is the reading
both Augustine and Cassiodorus expounded.

The Fathers divide over who prays the psalm. Augustine hears Christ, head and
body, praying in one voice. Cassiodorus declines to give the psalm to Christ,
because Christ is nowhere said to have been ``anxious'' and the ashes and tears
of the psalm do not befit his immaculate incarnation; he hears instead the
poor of Christ interceding for the whole world. Both hearings return in the
three readings.

\subsection{Gospel: the great commandment and the Son of David}

The Gospel is the last of the controversies in the Temple that run from Mt
21:23 to 22:46: the question of Jesus' authority, the tribute to Caesar
(22:15--22), the Sadducees' question about the resurrection (22:23--33), and
now the Pharisees. After it ``neither durst any man from that day forth ask him
any more questions'', and chapter 23 turns to the woes against the scribes and
Pharisees. When he has ended ``all these words'', Jesus tells his disciples,
``after two days shall be the pasch'' (26:1--2). The Missal's opening, \latin{Accesserunt ad Iesum pharisaei}, drops the
link with the silenced Sadducees and lets the pericope begin with the
Pharisees' approach.

The passage has two halves. In the first, a doctor of the law asks which is the
great commandment; Jesus answers with Deut 6:5, the love of God with the whole
heart, soul and mind, and adds unasked the second, Lev 19:18, ``like to this'',
on which with the first ``dependeth the whole law and the prophets''. In the
second, Jesus asks the Pharisees whose son the Christ is and confronts their
answer with Ps~109:1, ``The Lord said to my Lord''. The Fathers give different
reasons for the lawyer's question. Chrysostom says the Pharisees hoped that
Jesus would amend the first commandment ``for the sake of showing that Himself
too was God'' (\work{Homilies on Matthew} 71). Jerome says the question was
captious because every commandment of God is great, so that any answer could
be attacked. Aquinas reports both, and adds Augustine's reconciliation of
Matthew's ``tempting him'' with Mark's scribe who was ``not far from the
kingdom of God''. The first half is the anchor of the first reading and the
second half of the second.

The Gospel's place in the year has moved more than any other element of the
formulary. Fromage reports that it gave the day the name ``the Sunday of the
love of God'' only after the Gospel of the dropsical man and the wedding feast
(Luke 14) was moved eight days earlier. The medieval liturgical commentators
Rupert of Deutz and William Durandus, who expound this Sunday's office, both
read Luke~14 here. Durandus records that some churches read the present Gospel
on the following Sunday. In the 1962 Missal the pericopes that precede it in
Matthew come later in the year: the wedding feast (Mt 22:1--14) on the
Nineteenth Sunday and the tribute money (22:15--21) on the Twenty-second.

\subsection{Offertory: Daniel's prayer for the desolate sanctuary}

The Offertory comes from the ninth chapter of Daniel. In the first year of
Darius the Mede, Daniel understands from the book of Jeremias that the seventy
years of Jerusalem's desolation are nearly accomplished (9:1--2). He sets his
face to pray ``with fasting, and sackcloth, and ashes'' (9:3). He confesses the
sins of Israel and the justice of God: ``To thee, O Lord, justice: but to us
confusion of face'' (9:7); ``the Lord, our God, is just in all his works''
(9:14). He begs mercy for the desolate sanctuary, the city and the people
called by God's name (9:17--19). While he is still speaking, Gabriel comes to
him with the prophecy of the seventy weeks (9:20--27). The antiphon frames the
petitions of verses 17--19 with a narrative opening, \latin{Oravi Deum meum ego
Daniel, dicens}, formed from the chapter's own first-person voice, and it
condenses the petitions into three: hear thy servant, make thy face to shine
upon thy sanctuary, look with favour on this people.

Its \latin{illumina faciem tuam} is not the Vulgate's \latin{ostende}. It
agrees with the verb of Theodotion's Greek version of Daniel; Brenton's English
of the Greek Daniel reads ``cause thy face to shine on thy desolate
sanctuary''. The Hebrew, as the Revised Version translates it, has the same
sense. Which Latin version the chant followed has not been identified. Older books sang the antiphon with verses from the sequel:
\latin{Adhuc me loquente}, from Gabriel's coming, and a verse from Daniel~10
ending \latin{nam et Michael venit in adiutorium meum}. Fromage prints both.
He explains the choice of Daniel's prayer by the nearness of the feast of St
Michael (29 September): the ancient antiphonary he cites from Tommasi's edition
calls the following Sunday the first after Saint Michael, \latin{post Sancti
Angeli}. Durandus
likewise finds the angel fitting because the office is sung near the
dedication of St Michael. The 1962 Missal prints the antiphon alone. St Jerome's
commentary on Daniel is the principal patristic witness to the prayer itself;
it carries the third reading.

\subsection{Secret: stripped of past and future sins}

The Secret asks the divine majesty that ``these holy things which we do''
(\latin{haec sancta, quae gerimus}) may strip us, \latin{exuant}, as one strips
off a garment, both of past sins and of future ones. The 1861 English has
``cleanse us from all past offences, and from those we may hereafter be guilty
of''. The prayer stands with this Sunday's Collect in the Old Gelasian Mass of
Liber III, sect.~XIII, and in the Frankish Gelasians and the Gregorian
Supplement with both the Collect and the Postcommunion. The Council of Trent
teaches that Christ willed the Eucharist to be received \latin{tanquam
antidotum, quo liberemur a culpis quotidianis, et a peccatis mortalibus
praeservemur} (sess.~XIII, \work{Decretum de ss.~Eucharistia}, cap.~2), a
remedy for past faults and a safeguard against future ones. The same chapter
calls the sacrament \latin{symbolum unius illius corporis, cuius ipse caput
exsistit}, the sign of the one body of which Christ is head, bound together
\latin{arctissima fidei, spei et caritatis connexione}. The Secret's conclusion
is printed only as far as \latin{in unitate} and is said in full. The Missal
then directs the Preface of the Most Holy Trinity.

\subsection{Communion: vows before the God who is terrible to kings}

The Communion is taken from Ps~75, which the Vulgate titles ``for Asaph: a
canticle to the Assyrians''. The psalm sings of God known in Judea, whose
abode is in Sion and who there broke ``the powers of bows, the shield, the
sword, and the battle'' (vv.~2--4); of the day ``When God arose in judgment,
to save all the meek of the earth'' (v.~10); and, in the verses sung here, of
vows and gifts brought to the God ``who taketh away the spirit of princes''
(vv.~12--13). The Missal's \latin{apud omnes reges terrae} has an \latin{omnes}
that neither the Clementine Vulgate nor the psalters of Augustine and
Cassiodorus contain.

Augustine reads the vows and the ``spirit of princes'' morally, as the first
and third readings show. Cassiodorus reads \latin{omnes qui in circuitu eius affertis munera} of the
faithful who bring their gifts to the altar: \latin{In circuitu quippe res
agitur, cum munera fidelium sacratissimis altaribus offeruntur}. Bellarmine
gives the same reading, ``all you that are in the habit of approaching his
altars and offering your gifts upon them''. In both, the psalm's words are
heard at the place where the Communion is sung.

\subsection{Postcommunion: vices cured, eternal remedies}

The Postcommunion asks that by God's \latin{sanctificationes}, his sanctifying
gifts, ``our vices may be cured'' and ``eternal remedies'' may come to us. The
1861 English has ``an eternal remedy procured for us by these sacred
mysteries''. The prayer is older in Lent than at this Sunday. The first part of
Wilson's Gregorian Sacramentary gives it as the postcommunion of the Lenten
Ember Saturday and of Tuesday in Holy Week (there with \latin{sempiterna}), and
the 1962 Missal keeps it on both days. It reached the Sundays after Pentecost
with this Sunday's Collect and Secret in the Frankish Gelasians and the
Supplement.

The 1962 Missal also borrows its second clause for a later composition, the
Collect of St John of God (8~March), which asks, through that saint's merits,
\latin{ut igne caritátis tuæ vítia nostra curéntur, et remédia nobis ætérna
provéniant}: the clause stands there with \latin{igne caritátis tuæ} before it,
so that the cure is sought through the fire of God's charity. The Mass of St
John of God reads this Sunday's Gospel, Mt 22:34--46. That later Collect hears
the Postcommunion's cure as the work of charity, and sets it beside the Gospel
of the great commandment.

\subsection{The Preface of the Most Holy Trinity}

The Preface is not proper to this Sunday. The Missal appoints it for every
Sunday of the second class outside the Christmas and Easter seasons (p.~293,
no.~1082). It is nonetheless the confession this Mass makes before the Canon.
Its proper clause, \latin{Qui cum unigénito Fílio tuo, et Spíritu Sancto, unus
es Deus, unus es Dóminus}, confesses the one God with the Son and the Holy
Spirit, ``not in a singularity of one Person, but in a Trinity of one
substance'' (1861 English). It ends by adoring \latin{et in persónis
propríetas, et in esséntia únitas, et in maiestáte \dots\ æquálitas}. Its
``one God \dots\ one Lord'' stands close in wording to the Epistle's ``one
Lord \dots\ one God and Father''. That closeness returns in the second reading.

\subsection{Parts of different ages in one Mass}

No biblical author or ancient commentator wrote about this Mass as a finished
sequence, and, as these settings show, not every part reached the Sunday at
the same time. The three old chants already belonged together before
the present Gospel was appointed here; the Pauline lesson belongs to the
Missal's course through the Epistles; the prayers have their own sacramentary
history; and the Trinity Preface comes from the Sunday rubric rather than from
this formulary. Each element's own scriptural argument therefore comes first,
and what the liturgical sequence makes audible between the elements comes
after it.

That second hearing begins in three different places. One reading begins
with the commandment and asks how the Epistle's unity and the prayers of the
people sound when charity controls the hearing. One begins with Christ's
question and asks how the Epistle's confession, the Gradual's word and
spirit, and the appointed Preface sound when the identity of David's Lord
controls it. One begins with the older chant cycle and asks how every element
is received by a people that first confesses divine justice and then petitions
for mercy. Every part of the Mass enters each; in each only some texts
actually bear the claim, and at some places another reading explains more.

Each witness speaks of the passage he expounds, and the witnesses' agreements
carry the readings. They also disagree: Augustine and Cassiodorus hear
different speakers in the poor man, and Jerome and Chrysostom construe the
Epistle's final words differently, so that neither the psalm's speaker nor
the verse's construal is common ground. The patristic and saintly testimony
supplies the tested parts of each reading, and the editor's comparison of the
appointed texts the unity that joins them.
